% Ad Atticum 6.1 (ad-atticum-06-01)
% Source: https://raw.githubusercontent.com/PerseusDL/canonical-latinLit/master/data/phi0474/phi057/phi0474.phi057.perseus-lat1.xml
% Fetched by scripts/fetch_latin.py

% Dateline (Perseus): Scr. Laodiceae vi K. Mart. a. 704 (50).

\ciceroLetterOpener{CICERO ATTICO salutem}

\ciceroSection{1} accepi tuas litteras a. d. quintum Terminalia Laodiceae; quas legi libentissime plenissimas amoris, humanitatis, offici, diligentiae. iis igitur respondebo non χρύσεα χαλκείων (sic enim postulas) nec οἰκονομίαν meam instituam, sed ordinem conservabo tuum. recentissimas a Cybistris te meas litteras habere ais a. d. xi Kalendas Octobris datas et scire vis tuas ego quas acceperim. omnis fere quas commemoras, praeter eas quas scribis Lentuli pueris et Equotutico et Brundisio datas. qua re non οἴχεται tua industria quod vereris sed praeclare ponitur, si quidem id egisti ut ego delectarer. nam nulla re sum delectatus magis.

\ciceroSection{2} quod meam βαθύτητα in Appio tibi, liberalitatem etiam in Bruto probo, vehementer gaudeo; ac putaram paulo secus. Appius enim ad me ex itinere bis terve ὑπομεμψιμοίρουσ litteras miserat quod quaedam a se constituta rescinderem. ut si medicus, cum aegrotus alii medico traditus sit, irasci velit ei medico qui sibi successerit si quae ipse in curando constituerit mutet ille, sic Appius, cum ἐξ ἀφαιρέσεωσ provinciam curarit, sanguinem miserit, quicquid potuit detraxerit, mihi tradiderit enectam, προσανατρεφομένην eam a me non libenter videt sed modo suscenset, modo gratias agit. nihil enim a me fit cum ulla illius contumelia; tantum modo dissimilitudo meae rationis offendit hominem. quid enim potest esse tam dissimile quam illo imperante exhaustam esse sumptibus et iacturis provinciam, nobis eam obtinentibus nummum nullum esse erogatum nec privatim nec publice? quid dicam de illius praefectis, comitibus, legatis etiam? de rapinis, de libidinibus, de contumeliis? nunc autem domus me hercule nulla tanto consilio aut tanta disciplina gubernatur aut tam modesta est quam nostra tota provincia. haec non nulli amici Appi ridicule interpretantur qui me idcirco putent bene audire velle ut ille male audiat, et recte facere non meae laudis sed illius contumeliae causa. sin Appius, ut Bruti litterae quas ad te misit significabant, gratias nobis agit non moleste fero, sed tamen eo ipso die quo haec ante lucem scribebam, cogitabam eius multa inique constituta et acta tollere.

\ciceroSection{3} nunc venio ad Brutum quem ego omni studio te auctore sum complexus, quem etiam amare coeperam; sed ilico me revocavi, ne te offenderem. noli enim putare me quicquam maluisse quam ut mandatis satis facerem nec ulla de re plus laborasse. mandatorum autem mihi libellum dedit, isdemque de rebus tu mecum egeras. omnia sum diligentissime persecutus. primum ab Ariobarzane sic contendi ut talenta quae mihi pollicebatur illi daret. quoad mecum rex fuit, perbono loco res erat; post a Pompei procuratoribus sescentis premi coeptus est. Pompeius autem quom ob ceteras causas plus potest unus quam ceteri omnes, tum quod putatur ad bellum Parthicum esse venturus. ei tamen sic nunc solvitur, tricesimo quoque die talenta Attica xxxiii et hoc ex tributis. nec inde satis efficitur in usuram menstruam. sed Gnaeus noster clementer id fert; sorte caret, usura nec ea solida contentus est. alii neque solvit cuiquam nec potest solvere; nullum enim aerarium, nullum vectigal habet. Appi instituto tributa imperat. ea vix in faenus Pompei quod satis sit efficiunt. amici regis duo tresve perdivites sunt sed ii suum tam diligenter tenent quam ego aut tu. equidem non desino tamen per litteras rogare, suadere, accusare regem.

\ciceroSection{4} Deiotarus etiam mihi narravit se ad eum legatos misisse de re Bruti; eos sibi responsum rettulisse illum non habere. et me hercule ego ita iudico, nihil illo regno spoliatius, nihil rege egentius. itaque aut tutela cogito me abdicare aut ut pro Glabrione Scaevola faenus et impendium recusare. ego tamen quas per te Bruto promiseram praefecturas, M. Scaptio, L. Gavio, qui in regno rem Bruti procurabant, detuli; nec enim in provincia mea negotiabantur. tu autem meministi nos sic agere ut quot vellet praefecturas sumeret, dum ne negotiatori. itaque duas ei praeterea dederam. sed ii quibus petierat de provincia decesserant.

\ciceroSection{5} nunc cognosce de Salaminiis, quod video tibi etiam novum accidisse tamquam mihi. numquam enim ex illo audivi illam pecuniam esse suam; quin etiam libellum ipsius habeo, in quo est, Salaminii pecuniam debent M. Scaptio et P. Matinio, familiaribus meis. eos mihi commendat; adscribit etiam et quasi calcar admovet intercessisse se pro iis magnam pecuniam. confeceram ut solverent centesimis bienni ductis cum renovatione singulorum annorum. at Scaptius quaternas postulabat. metui , si impetrasset, ne tu ipse me amare desineres; nam ab edicto meo recessissem et civitatem in Catonis et in ipsius Bruti fide locatam meisque beneficiis ornatam funditus perdidissem.

\ciceroSection{6} atque hoc tempore ipso impingit mihi epistulam Scaptius Bruti rem illam suo periculo esse, quod nec mihi umquam Brutus dixerat nec tibi, etiam ut praefecturam Scaptio deferrem. id vero per te exceperamus ne negotiatori; quod si cuiquam, huic tamen non. fuerat enim praefectus Appio et quidem habuerat turmas equitum quibus inclusum in curia senatum Salamine obsederat, ut fame senatores quinque morerentur. itaque ego, quo die tetigi provinciam, cum mihi Cyprii legati Ephesum obviam venissent, litteras misi ut equites ex insula statim decederent. his de causis credo Scaptium iniquius de me aliquid ad Brutum scripsisse. sed tamen hoc sum animo. si Brutus putabit me quaternas centesimas oportuisse decernere, cum tota provincia singulas observarem itaque edixissem idque etiam acerbissimis faeneratoribus probaretur, si praefecturam negotiatori denegatam queretur, quod ego Torquato nostro in tuo Laenio, Pompeio ipsi in Sex. Statio negavi et iis probavi, si equites deductos moleste feret, accipiam equidem dolorem mihi illum irasci sed multo maiorem non esse eum talem qualem putassem.

\ciceroSection{7} illud quidem fatebitur Scaptius, me ius dicente sibi omnem pecuniam ex edicto meo auferendi potestatem fuisse. addo etiam illud quod vereor tibi ipsi ut probem. consistere usura debuit quae erat in edicto meo. deponere volebant: impetravi a Salaminiis ut silerent. veniam illi quidem mihi dederunt, sed quid iis fiet, si huc Paulus venerit? sed totum hoc Bruto dedi; qui de me ad te humanissimas litteras scripsit, ad me autem, etiam cum rogat aliquid, contumaciter, adroganter, ἀκοινονοήτωσ solet scribere. tu autem velim ad eum scribas de his rebus, ut sciam quo modo haec accipiat; facies enim me certiorem. atque haec superioribus litteris diligenter ad te per scripseram sed plane te intellegere volui mihi non excidisse illud quod tu ad me quibusdam litteris scripsisses, si nihil aliud de hac provincia nisi illius benevolentiam deportassem, mihi id satis esse. sit sane, quoniam ita tu vis, sed tamen cum eo credo quod sine peccato meo fiat. igitur meo decreto soluta res Scaptio stat. quam id rectum sit tu iudicabis; ne ad Catonem quidem provocabo.

\ciceroSection{8} sed noli me putare ἐγκελεύσματα illa tua abiecisse quae mihi in visceribus haerent. flens mihi meam famam commendasti; quae epistula tua est in qua non eius mentionem facias? itaque irascatur qui volet; patiar. τὸ γὰρ εὖ praesertim cum sex libris tamquam praedibus me ipse obstrinxerim, quos tibi tam valde probari gaudeo. e quibus unum ἱστορικὸν requiris de Cn. Flavio, Anni filio. ille vero ante decemviros non fuit quippe qui aedilis curulis fuerit, qui magistratus multis annis post decemviros institutus est. quid ergo profecit quod protulit fastos? occultatam putant quodam tempore istam tabulam, ut dies agendi peterentur a paucis. nec vero pauci sunt auctores Cn. Flavium scribam fastos protulisse actionesque composuisse, ne me hoc vel potius Africanum (is enim loquitur) commentum putes. οὐκ ἔλαθέ σε illud de gestu histrionis. tu sceleste suspicaris, ego ἀφελῶσ scripsi. de me imperatore scribis te ex Philotimi litteris cognosse; sed credo te, iam in Epiro cum esses, binas meas de omnibus rebus accepisse, unas a Pindenisso capto, alteras Laodicea, utrasque tuis pueris datas. quibus de rebus propter casum navigandi per binos tabellarios misi Romam publice litteras.

\ciceroSection{10} de Tullia mea tibi adsentior scripsique ad eam et ad Terentiam mihi placere. tu enim ad me iam ante scripseras, ac vellem te in tuum veterem gregem rettulisses. correcta vero epistula Memmiana nihil negoti fuit; multo enim malo hunc a Pontidia quam illum a Servilia. qua re adiunges Saufeium nostrum, hominem semper amantem mei, nunc credo eo magis quod debet etiam fratris Appi amorem erga me cum reliqua hereditate crevisse; qui declaravit quanti me faceret cum saepe tum in Bursa. ne tu me sollicitudine magna liberaris.

\ciceroSection{11} Furni exceptio mihi non placet; nec enim ego ullum aliud tempus timeo nisi quod ille solum excipit. sed scriberem ad te de hoc plura, si Romae esses. in Pompeio te spem omnem oti ponere non miror. ita res est removendumque censeo illud dissimulantem. sed enim οἰκονομία si perturbatior est, tibi adsignato. te enim sequor σχεδιάζοντα .

\ciceroSection{12} Cicerones pueri amant inter se, exercentur, sed discunt, alter, uti dixit Isocrates in Ephoro et Theopompo, frenis eget, alter calcaribus. Quinto togam puram Liberalibus cogitabam dare; mandavit enim pater. ea sic observabo quasi intercalatum non sit. Dionysius mihi quidem in amoribus est; pueri autem aiunt eum furenter irasci; sed homo nec doctior nec sanctior fieri potest nec tui meique amantior.

\ciceroSection{13} Thermum, Silium vere audis laudari. valde honeste se gerunt. adde M. Nonium, Bibulum, me, si voles. iam Scrofa vellem haberet ubi posset; est enim lautum negotium. ceteri infirmant πολίτευμα Catonis. Hortensio quod causam meam commendas valde gratum. de Amiano spei nihil putat esse Dionysius. Terenti nullum vestigium adgnovi. Moeragenes certe periit. feci iter per eius possessionem in qua animal reliquum nullum est. haec non noram tum, cum Democrito tuo cum locutus sum. Rhosica vasa mandavi. sed heus tu! quid cogitas? in felicatis lancibus et splendidissimis canistris holusculis nos soles pascere; quid te in vasis fictilibus appositurum putem? κέρασ Phemio mandatum est; reperietur, modo aliquid illo dignum canat.

\ciceroSection{14} Parthicum bellum impendet. Cassius ineptas litteras misit, necdum Bibuli erant adlatae. quibus recitatis puto fore ut aliquando commoveatur senatus. equidem sum in magna animi perturbatione. si , ut opto, non prorogatur nostrum negotium, habeo Iunium et Quintilem in metu. esto ; duos quidem mensis sustinebit Bibulus. quid illo fiet quem reliquero, praesertim si fratrem? quid me autem, si non tam cito decedo? Magna turba est. mihi tamen cum Deiotaro convenit ut ille in meis castris esset cum suis copiis omnibus. habet autem cohortis quadringenarias nostra armatura XXX , equitum cIↃ cIↃ . erit ad sustentandum quoad Pompeius veniat; qui litteris quas ad me mittit significat suum negotium illud fore. hiemant in nostra provincia Parthi; exspectatur ipse Orodes. quid quaeris? aliquantum est negoti. de Bibuli edicto nihil novi praeter illam exceptionem de qua tu ad me scripseras nimis gravi praeiudicio in ordinem nostrum. ego tamen habeo ἰσοδυναμοῦσαν sed tectiorem ex Q. Muci P. L edicto Asiatico, extra quam si ita negotium gestum est ut eo stari non oporteat ex fide bona, multaque sum secutus Scaevolae, in iis illud in quo sibi libertatem censent Graeci datam, ut Graeci inter se disceptent suis legibus. breve autem edictum est propter hanc meam διαίρεσιν quod duobus generibus edicendum putavi. quorum unum est provinciale in quo est de rationibus civitatum, de aere alieno, de usura, de syngraphis, in eodem omnia de publicanis; alterum, quod sine edicto satis commode transigi non potest, de hereditatum possessionibus, de bonis possidendis, vendendis, magistris faciendis, quae ex edicto et postulari et fieri solent. tertium de reliquo iure dicundo ἄγραφον reliqui. dixi me de eo genere mea decreta ad edicta urbana accommodaturum. itaque curo et satis facio adhuc omnibus. Graeci vero exsultant quod peregrinis iudicibus utuntur. nugatoribus quidem inquies. quid refert? tamen se αὐτονομίαν adeptos putant. vestri enim credo gravis habent Turpionem sutorium et Vettium mancipem.

\ciceroSection{16} de publicanis quid agam videris quaerere. habeo in deliciis, obsequor, verbis laudo, orno; efficio ne cui molesti sint. τὸ , usuras eorum quas pactionibus adscripserant servavit etiam Servilius. ego sic. diem statuo satis laxam, quam ante si solverint, dico me centesimas ducturum; si non solverint, ex pactione. itaque et Graeci solvunt tolerabili faenore et publicanis res est gratissima, si illa iam habent pleno modio, verborum honorem, invitationem crebram. quid plura? sunt omnes ita mihi familiares ut se quisque maxime putet. sed tamen μηδὲν αὐτοῖσ —scis reliqua.

\ciceroSection{17} de statua Africani ( ὢ πραγμάτων ! sed me id ipsum delectavit in tuis litteris) ain tu? Scipio hic Metellus proavum suum nescit censorem non fuisse? atqui nihil habuit aliud inscriptum nisi cos ea statua quae ad Opis †per te† posita in excelso est. in illa autem quae est ad Πολυκλέουσ Herculem inscriptum est consul ; quam esse eiusdem status, amictus, anulus, imago ipsa declarat. at me hercule ego, cum in turma inauratarum equestrium quas hic Metellus in Capitolio posuit animadvertissem in Serapionis subscriptione Africani imaginem, erratum fabrile putavi, nunc video Metelli.

\ciceroSection{18} o ἀνιστορησίαν turpem! nam illud de Flavio et fastis, si secus est, commune erratum est et tu belle ἠπόρησασ et nos publicam prope opinionem secuti sumus, ut multa apud Graecos. quis enim non dixit Εὔπολιν τὸν τῆσ ἀρχαίασ ab Alcibiade navigante in Siciliam deiectum esse in mare? redarguit Eratosthenes; adfert enim quas ille post id tempus fabulas docuerit. num idcirco Duris Samius, homo in historia diligens, quod cum multis erravit, inridetur? quis Zaleucum leges Locris scripsisse non dixit? num igitur iacet Theophrastus si id a Timaeo tuo familiari reprensum est? sed nescire proavum suum censorem non fuisse turpe est, praesertim cum post eum consulem nemo Cornelius illo vivo censor fuerit.

\ciceroSection{19} quod de Philotimo et de solutione HS X_X_DC scribis, Philotimum circiter Kal. Ianuarias in Chersonesum audio venisse. at mi ab eo nihil adhuc. reliqua mea Camillus scribit se accepisse. ea quae sint nescio et aveo scire. verum haec posterius et coram fortasse commodius.

\ciceroSection{20} illud me, mi Attice, in extrema fere parte epistulae commovit; scribis enim sic, τί λοιπόν ; deinde me obsecras amantissime ne obliviscar vigilare et ut animadvertam quae fiant. num quid de quo inaudisti? etsi nihil eius modi est. πολλοῦ γε . nec enim me fefellisset nec fallet. sed ista admonitio tua tam accurata nescio quid mihi significare visa est.

\ciceroSection{21} de M. Octavio iterum iam tibi rescribo te illi probe respondisse; paulo vellem fidentius. nam Caelius libertum ad me misit et litteras accurate scriptas et de pantheris et† a civitatibus. rescripsi alterum me moleste ferre, si ego in tenebris laterem nec audiretur Romae nullum in mea provincia nummum nisi in aes alienum erogari, docuique nec mihi conciliare pecuniam licere nec illi capere monuique eum quem plane diligo ut cum alios accusasset cautius viveret; illud autem alterum alienum esse existimatione mea, Cibyratas imperio meo publice venari.

\ciceroSection{22} Lepta tua epistula gaudio exsultat; etenim scripta belle est meque apud eum magna in gratia posuit. filiola tua gratum mihi fecit quod tibi diligenter mandavit ut mihi salutem adscriberes, gratum etiam Pilia, sed illa officiosius quod mihi quem iam pridem numquam vidit. igitur tu quoque salutem utrique adscribito. litterarum datarum dies pr. Kal. Ianuar. suavem habuit recordationem clarissimi iuris iurandi quod ego non eram oblitus. Magnus enim praetextatus illo die fui. habes ad omnia. non , ut postulasti, χρύσεα χαλκείων sed paria paribus respondimus.

\ciceroSection{23} ecce autem alia pusilla epistula quam non relinquam ἀναντιφώνητον . bene me hercule †potuit Lucceius Tusculanum, nisi forte (solet enim) cum suo tibicine†. et velim scire qui sit eius status. Lentulum quidem nostrum omnia praeter Tusculanum proscripsisse audio. cupio hos expeditos videre, cupio etiam Sestium, adde sis Caelium; in quibus omnibus est αἴδεσθεν μὲν de Memmio restituendo ut Curio cogitet te audisse puto. de Egnati Sidicini nomine nec nulla nec magna spe sumus. Pinarium quem mihi commendas diligentissime Deiotarus curat graviter aegrum. respondi etiam minori.

\ciceroSection{24} tu velim dum ero Laodiceae, id est ad Idus Maias, quam saepissime mecum per litteras colloquare et cum Athenas veneris (iam enim sciemus de rebus urbanis, de provinciis, quae omnia in mensem Martium sunt conlata), utique ad me tabellarios mittas.

\ciceroSection{25} et heus tu! †genuarios† a Caesare per Herodem talenta Attica L extorsistis? in quo, ut audio, magnum odium Pompei suscepistis. putat enim suos nummos vos comedisse, Caesarem in Nemore aedificando diligentiorem fore. haec ego ex P. Vedio, magno nebulone sed Pompei tamen familiari, audivi. hic Vedius mihi obviam venit cum duobus essedis et raeda equis iuncta et lectica et familia magna pro qua, si Curio legem pertulerit, HS centenos pendat necesse est. erat praeterea cynocephalus in essedo nec deerant onagri. numquam vidi hominem nequiorem. sed extremum audi. deversatus est Laodiceae apud Pompeium Vindullum. lbi sua deposuit cum ad me profectus est. moritur interim Vindullus; quae res ad Magnum Pompeium pertinere putabatur. C. Vennonius domum Vindulli venit. cum omnia obsignaret, in Vedianas res incidit. in his inventae sunt quinque imagunculae matronarum in quibus una sororis amici tui hominis bruti qui hoc utatur et illius lepidi qui haec tam neglegenter ferat. haec te volui παριστορῆσαι . sumus enim ambo belle curiosi.

\ciceroSection{26} unum etiam velim cogites. audio Appium πρόπυλον Eleusine facere. num inepti fuerimus si nos quoque Academiae fecerimus? puto inquies. ergo id ipsum scribes ad me. equidem valde ipsas Athenas amo. volo esse aliquod monumentum; odi falsas inscriptiones statuarum alienarum. sed ut tibi placebit, faciesque me in quem diem Romana incidant mysteria certiorem et quo modo hiemaris. cura ut valeas. post Leuctricam pugnam die septingentesimo sexagesimo quinto.
